\section{When the optimizer reaches the training optimum}\label{app:convex}
We froze the hidden rectifiers of a shallow scalar network and fitted only its output coefficients, including the direct input slope and constant offset. For each target (square, sine, positive self-power and hyperbolic tangent), we used seven hidden widths \(0,1,3,7,15,31,63\) and three random starts. The frozen features were rectifiers with Xavier slopes and zero biases, rectifiers with equally spaced central knots, or rectifiers with equally spaced knots throughout the deployment range. A knot is the input where a rectifier changes slope. Each arrangement used the same initial output weights for a given width and start. Fits used 1,024 independent labeled examples sampled uniformly from either the central range or, in a separate full-information control, the deployment range. Independent test draws comprised 4,096 central examples and 4,096 per applicable outer side.

For observation \(i\), let \(x_i\) be its input, \(y_i\) its label and \(s_i=\max(1,|y_i|)\) its loss scale. Stack its frozen feature values, \(x_i\) and the constant 1 into a row and divide it by \(s_i\). Let \(M\) be the resulting design matrix, and let column \(b\) have entries \(y_i/s_i\). Write \(\theta\) for the output coefficients, \(\theta_0\) for their initial values, and \({}^\top\) for matrix transpose. The training residual vector is \(M\theta-b\). The Moore--Penrose pseudoinverse \(M^+\) maps any residual vector to its minimum-length least-squares coefficient vector. The nearest global minimizer in Euclidean distance to \(\theta_0\) is
\[
             \theta_*=\theta_0-M^+(M\theta_0-b).
\]
At \(\theta_*\), the residual is perpendicular to every design column, which certifies the minimum in real arithmetic.

The singular value decomposition (SVD) expresses \(M\) as orthogonal input and output directions with nonnegative gains called singular values. Let \(\sigma_j\) be gain \(j\), \(\sigma_{\max}\) the largest, \(m,n\) the numbers of rows and columns, and \(\epsilon=2^{-52}\) the spacing immediately above 1 in binary64. We keep directions with \(\sigma_j>\sigma_{\max}\max(m,n)\epsilon\) when computing the pseudoinverse. The release records the largest absolute entry of \(M^\top(M\theta-b)\) divided by \(m\). This is zero when the normal equations hold exactly. In exact arithmetic, one gradient update on half the summed squared residual, preconditioned by \((M^\top M)^+\), gives \(\theta_*\). We also compute an idealized 5,000-step trajectory of full-batch gradient descent (GD) on that summed loss with step size \(1/\sigma_{\max}^2\). In each retained singular direction its residual component is multiplied by \((1-\sigma_j^2/\sigma_{\max}^2)^{5000}\). Evaluating this formula avoids 5,000 separate matrix multiplications. The reported training loss divides the summed loss by the number of observations.

\paragraph{Sixty-three centrally placed hinges.}
The table reports the median of three starts. Training loss uses the scaled labels above. Outer mean scaled error averages \( |\hat y-y|/\max(1,|y|) \) over held-out scalar predictions \(\hat y\) with references \(y\). Outer root mean squared error (RMSE) is in original target units. GD and SVD are two solvers for the same frozen-feature fit. The preconditioned gradient calculation agrees with SVD to the digits shown.
\begin{center}
\begin{minipage}{\linewidth}
\centering
\makeatletter\def\@captype{table}\makeatother
\caption{Central loss, mean scaled outer error and unscaled outer RMSE on frozen features. Each result is the median of three starts for 63 central hinges. GD uses 5,000 spectral steps and SVD uses the numerical least-squares solution.}\label{tab:convex_rmse}
\small\setlength{\tabcolsep}{4pt}
\begin{tabular}{@{}lrrrrrr@{}}\toprule
& \multicolumn{2}{c}{Central training loss} & \multicolumn{2}{c}{Outer scaled error} & \multicolumn{2}{c}{Outer RMSE}\\
Target & GD & SVD & GD & SVD & GD & SVD\\\midrule
Square & $.00149$ & $8.68\times 10^{-8}$ & .506 & .275 & 34.0 & 22.3\\
Sine & $.00113$ & $2.59\times 10^{-8}$ & 2.06 & 3.27 & 3.05 & 4.09\\
Positive self-power & $.00159$ & $2.69\times 10^{-7}$ & .932 & .882 & $1.43\times 10^9$ & $1.43\times 10^9$\\
Tanh & $.000920$ & $9.32\times 10^{-9}$ & .952 & .0333 & 1.20 & .0396\\\bottomrule
\end{tabular}
\end{minipage}
\end{center}
The globally optimized central sine fit extends its last observed slope into the shell, producing higher outer RMSE than finite-step GD. With knots spread through the full envelope, central sine labels give a numerical design rank of 21 out of 65 directions. Permitting labels throughout the envelope raises the rank to 65 and lowers mean outer scaled error from 5.43 to .00194 with the same optimizer. In the zero-bias Xavier arrangement all initial hinges meet at zero, so the model has at most two affine rays and design rank at most three. End-to-end BP can move those hinges.
